\subsection{The remaining ordinary five-state shape at order thirteen}

\begin{proposition}[Proved: three columns and four outside elements]
\label{CP2:A-three-by-three-four-outside-five-state}
Let $M$ have rectangular minimal ideal $K$ of shape $3\times3$
and at most four outside elements. An ordinary right-act cone
over $M$ with at most five states cannot have type $\CP^2$.
\end{proposition}
\begin{proof}
This extends the core-orbit argument of
Lemma~\ref{CP2:A-three-by-three-five-state}; we give the full
argument with the new bound. Write $S$ for the nonempty state
set and $q=|M\setminus K|\leq4$. Every core orbit in $SK$ is
a quotient of the three-column right ideal $eM$. It therefore
has one, two or three states, with a surjective column map.
Each orbit is individually invariant under $M$, because
\[
 s(i,j)m=s(i,\tau_m(j)).
\]
A one-state core orbit is consequently a global fixed state.
It gives a section of the action-bar projection to $BM$, so
all second classes of the act cone are relative to the monoid
vertex and their products vanish. We may exclude this case.

If there is exactly one core orbit, every band element acts
constantly on $S$. Indeed its value on any state belongs to
that orbit and is the orbit's value in the prescribed column.
The actual constant-action compression of
Theorem~\ref{CP2:A-constant-act-compression} gives an equivalent
monoid with the same $q$ outside elements and minimal band
$\{0,1,2\}\times S$. Its order is
\[
 q+3|S|\leq4+15=19<23.
\]
This contradicts the complete monoid minimum
Theorem~\ref{A:monoid-minimum23-classification}. The comparison
here is an actual homotopy equivalence, not merely a rational
second-cohomology comparison.

Suppose there are at least two core orbits. Since no orbit is
a singleton and $|S|\leq5$, there are exactly two. If their
union is $S$, its two invariant pieces give a disconnected
action category. The cone's first rational homology is then
nonzero, whereas a projective plane has zero first homology.
Thus the only remaining case consists of two two-point core
orbits and a single transient state $x$.

Let $\pi_0,\pi_1$ be the two surjections from the three old
columns onto the two core orbits. If their two-part partitions
differ, the bipartite graph with the four parts as vertices
and the three columns as edges is connected. To see this
directly, a disconnected graph without isolated vertices on
two vertices on either side would have two separate paired
components, forcing the partitions to agree. In the actual
constant-action quotient the relations
$\pi_0(j)\sim\pi_1(j)$ therefore identify all four core states.
Their common image is a global fixed state. The quotient cone
has zero second products, and its actual pullback surjects on
$H^2$ by Lemma~\ref{CP2:A-constant-quotient-H2}. Hence the original
cone has zero second products as well.

If the partitions agree, write $\pi$ for their common quotient.
The action of each outside element preserves either core orbit,
so the identity $s_jm=s_{\tau_m(j)}$ shows that $\tau_m$
preserves equality of $\pi$-values. Collapsing just these equal
columns of $K$, while keeping every outside element distinct,
is therefore an actual monoid congruence. Its quotient $M'$
has minimal band $3\times2$ and the same $q\leq4$ outside
elements. The entire action on $S$ factors through $M'$:
this holds on the core orbits by definition; for $x$ and a
fixed row $i$, the map $j\mapsto x(i,j)$ is the column map
of one of those same two core orbits. The actual functor of
Lemma~\ref{CP2:A-column-quotient-act-H2} is an isomorphism on
rational second cohomology. The target has zero second products
unconditionally by Proposition~\ref{CP2:A-three-row-five-outside-acts}.
Naturality gives the contradiction. No higher boundedness is
transferred through this column quotient.
\end{proof}

\begin{proposition}[Proved: a single remaining ordinary order-thirteen shape]
\label{CP2:A-thirteen-five-state-single-shape}
If an ordinary right-act cone over a monoid of order at most
thirteen with at most five states has type $\CP^2$, then its
source has order exactly thirteen, its minimal ideal is a
$4\times2$ rectangular band, and there are exactly five outside
elements. In particular the proposition is a shape reduction;
it does not exclude this remaining shape or the entire
nineteen-arrow budget $(13,1,5)$.
\end{proposition}
\begin{proof}
Theorem~\ref{CP2:A-bounded-twelve-acts} covers source orders through
twelve. A possible nonzero second square requires nonzero
$H^2(BM;\mathbb Q)$ and at least three minimal rows, by the
relative-product and two-row ordinary-act obstructions. It
also requires at least two minimal columns. Write the minimal
ideal as $I\times G\times J$, with $r=|I|\geq3$ and
$c=|J|\geq2$. Its outside identity gives
$rc|G|+1\leq13$.

If $G$ is nontrivial, the only possibility is
$r=3,c=2,|G|=2$ and one outside element. The actual minimal-group
collapse of Theorem~\ref{CP2:A-minimal-group-collapse-collage-H2},
with trivial opposite corner, preserves rational second
cohomology. Its target source is the pure $3\times2$ band with
identity. The contracted algebra is $\mathbb Q$, so its row
module is projective; Lemma~\ref{CP2:A-act-row-pd-one} kills
every second product on every induced act cone. This is an
unconditional target obstruction and uses no transfer of
boundedness through group collapse.

The minimal ideal is therefore a band. The complete list of
possible shapes, with the corresponding upper outside bounds,
is
\[
\begin{array}{c|rrrrrrr}
 (r,c)&(3,2)&(4,2)&(5,2)&(6,2)&(3,3)&(3,4)&(4,3)\\\hline
 q\leq&7&5&3&1&4&1&1.
\end{array}
\]
The $3\times2$ case is
Theorem~\ref{CP2:A-three-row-seven-outside-acts}. The
$5\times2$ and $6\times2$ cases are covered by
Lemma~\ref{CP2:A-two-column-four-outside-bounded}, applied to
the actual bounded cone. The $3\times3$ case is the preceding
proposition. In the $3\times4$ and $4\times3$ cases the single
outside element is the identity, so the contracted algebra is
again $\mathbb Q$ and the row module is projective. The same
unconditional projective-dimension lemma applies.
Only $4\times2$ with five outside elements remains; a smaller
outside count would put the source at order at most twelve.
\end{proof}

\begin{remark*}[Actual versus rational scope]
The one-core branch of the first proposition uses the actual
monoid minimum. It is not a proof excluding every rational
projective-plane ring on an arbitrary five-state quotient.
Likewise the second proposition leaves the order-two-opposite
budget $(13,2,4)$ separate. No assertion about either open
whole budget follows from the present shape reduction alone.
\end{remark*}
